%! Displaying and Analyzing Categorical Data — Unit 3, Lesson 3.1 — Lesson Plan
% GRADUAL-RELEASE plan, Unit 2-on timing 5 / 45 / 5 / 5: warm-up / guided notes and practice /
% spoken debrief / close and ASSIGN the homework (not started in class). One Main Ideas / Notes
% table: I do / we do in rows 1-4 (problems 1-11, ~25 min) -> we do the Guided Practice row
% (problems 12-15, ~20 min). THE HOMEWORK IS THE INDIVIDUAL PRACTICE. Teacher-facing. Breakable
% boxes are `skillbox` with a \boxguard ahead; the two tabularx boxes are `fixedskillbox`.
% Unit 3 condensed 2026-11-09: this lesson folds in the old unit launch and old Lesson 3.1.
\documentclass[10pt]{article}
\usepackage{saar-article}
\usepackage{saar-boxes}
\usepackage{graphicx}   % -article does NOT load graphicx
\graphicspath{{images/}}
\newcommand{\TallMath}[1]{$\displaystyle #1\rule[-1.4em]{0pt}{3.2em}$}
\newcommand{\UnitNumberName}{Unit 3: Categorical Data and Probability \quad}
\newcommand{\LessonNumberName}{Lesson 3.1: Displaying and Analyzing Categorical Data}

\begin{document}
\begin{center}
    {\Large\bfseries \CourseName \\
    {\small\bfseries \UnitNumberName \LessonNumberName}}
\end{center}

\begin{tcolorbox}[colback=frost, colframe=cerulean, boxrule=0.9pt, arc=2mm,
    left=2mm, right=2mm, top=1.5mm, bottom=1.5mm]
    \textbf{Primary Objective:} Students will read bar graphs of one categorical variable as
    frequencies and relative frequencies; compare groups with side-by-side and segmented bar
    graphs; generate a two-way table from a summary and read its joint counts and marginal
    frequencies; calculate and interpret marginal and conditional relative frequencies in
    context; and justify, with conditional percents, which of two groups of different sizes is
    more likely to fall in a category.
    \\[2pt]
    \textbf{Standards:} PS.DS.4a, PS.DS.4b, PS.DS.4c, PS.DS.4d \quad$\cdot$\quad the first lesson
    of Unit~3 (there is no separate launch lesson: the warm-up and the hook open the unit)
    \\[2pt]
    \textbf{Lesson model:} \emph{Gradual release --- I do, we do, you do.} There is no group
    activity. The notes name the vocabulary as they teach it and deliver the instruction
    \textbf{row by row in one Main Ideas / Notes table} (four rows, problems 1--11: you read the
    definition, mark up the display, and work the first problem of each grid; the class works the
    rest); the Guided Practice row (problems 12--15) is worked together with students holding the
    pen. \textbf{The homework is the individual practice} --- assigned at Close \& Assign and
    done outside class, not started in class. The debrief is a \emph{spoken} cold check and one
    board item; there is no exit ticket.
    \\[2pt]
    \textbf{Conventions on the page:} percents are exact except where a problem says about;
    charts and the library's table are pre-made spreadsheet output to read, so nobody draws a
    graph.
\end{tcolorbox}

\renewcommand{\arraystretch}{1.4}
\boxguard
\begin{skillbox}[Learning Targets \& Key Understandings]{goldbox}
\begin{tabularx}{\linewidth}{@{}X|X@{}}
\textbf{Learning targets (student language)} & \textbf{Key understandings (the ``why'')} \\[2pt]
\hline
\vspace{-6pt}
\begin{itemize}[leftmargin=1.1em, nosep, after=\vspace{2pt}]
  \item I can read a bar graph of one categorical variable in context, as a frequency or a
        \textbf{relative frequency}. \hfill \textit{PS.DS.4a}
  \item I can compare groups with \textbf{side-by-side} and \textbf{segmented} bar graphs.
        \hfill \textit{PS.DS.4b}
  \item I can build a \textbf{two-way table} from a summary and read its joint counts and
        \textbf{marginal frequencies}. \hfill \textit{PS.DS.4c}
  \item I can calculate and interpret marginal and \textbf{conditional relative frequencies},
        and name the whole each percent is taken of. \hfill \textit{PS.DS.4d}
  \item I can compare groups of different sizes with conditional percents, not raw counts, and
        explain why. \hfill \textit{PS.DS.4d}
\end{itemize}
&
\vspace{-6pt}
\begin{itemize}[leftmargin=1.1em, nosep, after=\vspace{2pt}]
  \item \textbf{Target misconception:} comparing raw \emph{counts} across groups of different
        sizes --- ``more ninth graders ride, so ninth graders are more likely to ride.''
        \textbf{The rule that replaces it:} counts say how many; conditional percents say how
        likely. To compare groups of different sizes, compare conditional percents.
  \item \textbf{Second misconception:} reading the wrong denominator (a row percent for a column
        percent). $54/94$ is a percent of bus riders, not of ninth graders. \textbf{The test:} in
        ``Of the [group], what percent \ldots?'' the words after \emph{of} name the whole.
  \item A two-way table summarizes two categorical variables at once, and its Totals are the
        wholes that percents are built on. Graphs and tables are two views of one set of counts
        (PS.DS.4).
  \item A relative frequency is the idea Lesson~3.2 turns into a probability; a conditional
        relative frequency is the idea Lesson~3.4 names conditional probability.
\end{itemize}
\end{tabularx}
\end{skillbox}

\boxguard[24]
\begin{skillbox}[{Vocabulary, Concepts, \& Theorems --- taught directly in the guided notes}]{greenbox}
\small
\renewcommand{\arraystretch}{1.35}
\begin{tabularx}{\linewidth}{@{} >{\bfseries}p{3.4cm} X @{}}
Relative frequency & A category's count divided by the total, as a fraction, decimal, or
       percent. \\
Side-by-side / segmented bars & Side-by-side bars stand the groups next to each other within each
       category; segmented bars stack each group's categories into one bar. \\
Two-way table & A table that records two categorical variables at once, one in the rows and one
       in the columns. Each inside cell is a joint count. \\
Marginal frequency & A count in the Total row or Total column: one variable alone. A
       \emph{marginal relative frequency} divides that Total by the whole table. \\
Conditional relative frequency & An inside count divided by the total of one group (one row or
       one column): the percent \emph{of that group}. \\
\end{tabularx}
\end{skillbox}

% fixedskillbox never splits, so the phase table stays intact on one page.
% THE PHASE TABLE IS A CONTRACT: 5 / 45 / 5 / 5 — I do ~25 of the 45, Guided Practice ~20.
\begin{fixedskillbox}[Lesson at a Glance --- \MeetingLength]{frost}
\renewcommand{\arraystretch}{1.25}
\begin{tabularx}{\linewidth}{@{}l c X X@{}}
\textbf{Phase} & \textbf{Min} & \textbf{Students} & \textbf{Teacher} \\
\hline
Warm-Up & 5 & Percents of a number, fractions as percents, a simple table --- silent & Debrief
  item~2's last question aloud; write \emph{larger top number, smaller percent} on the board and
  leave it up all period \\
Guided Notes \& Practice & 45 & Fill the vocabulary box as each term is named; work rows 1--4
  (problems 1--11) with the pen in their hand; then the Guided Practice row, \emph{Streaming
  Choices} (12--15), together & \textbf{I do / we do, row by row}: read the definition, mark up
  the display, work problems 1, 3, 6, 8; the class works the rest (25) $\to$ \textbf{we do}
  Guided Practice 12--15 (20) \\
Debrief & 5 & Pens down --- answer aloud, nothing to fill in & Put problem~15 back on the board
  ($40\%$ vs.\ $45\%$); ask the cold check \\
Close \& Assign & 5 & Write down the homework, the due date, and item~5 as the first item &
  Announce the homework and its form (DeltaMath + packet items 2, 5, 6), preview Lesson~3.2 \\
\end{tabularx}
\end{fixedskillbox}

\begin{fixedskillbox}[Warm-Up --- Activate Prior Knowledge (5 min)]{frost}
\begin{tabularx}{\linewidth}{@{}X X@{}}
\begin{minipage}[t]{\linewidth}\vspace{0pt}
    \textbf{The three items and what each seeds}
    \begin{itemize}[leftmargin=1.1em, itemsep=1pt, topsep=2pt]
      \item \textbf{1.} $47\%$ of $200$, $45\%$ of $120$, $50\%$ of $80$ give $94$, $54$, $40$.
            \textbf{Spirals} to percent work from earlier courses --- and those are row~1's bus
            bar and row~2's two bus bars.
      \item \textbf{2.} $94/200$, $54/120$, $40/80$ are $47\%$, $45\%$, $50\%$. \textbf{Seeds:}
            problems 8 and~9, and the crux: $54/120$ has the larger top number and the smaller
            percent.
      \item \textbf{3.} The one-variable table: $80$ by car, $26/200 = 13\%$ walk.
            \textbf{Seeds:} row~1's two bar graphs, which are this table drawn.
    \end{itemize}
\end{minipage}
&
\begin{minipage}[t]{\linewidth}\vspace{0pt}
    \textbf{Running it}
    \begin{itemize}[leftmargin=1.1em, itemsep=1pt, topsep=2pt]
      \item \textbf{Say what Unit~3 is:} data sorted into groups, and the chances those data
            give us. That sentence and the hook are the whole unit launch.
      \item \textbf{Debrief item~2's last question aloud} and write \emph{larger top number,
            smaller percent} on the board. Leave it up all period; row~4 points back to it.
      \item Item~2 is the tell. A student who says the larger top number must be the larger
            percent has walked in with today's misconception --- note who, and cold-call them
            at problem~9.
    \end{itemize}
\end{minipage}
\end{tabularx}
\end{fixedskillbox}

\boxguard
\begin{skillbox}[Hook --- before anyone sits down]{frost}
The hook is on \textbf{page 1 of the notes}, under the vocabulary box, and on the board:
\textit{``$54$ ninth graders ride the bus and only $40$ twelfth graders do, so ninth graders are
more likely to ride the bus.''} \textbf{Ask: do you agree?} Students circle \emph{Agree},
\emph{Disagree}, or \emph{Can't tell yet} and write one line of reason \emph{before} anyone
speaks; then take the hand vote and write the split on the board. Most rooms agree. \textbf{Do
not settle it.} Say only that the answer is problem~9, in the \emph{Which Whole?} row, and that
they will vote again there. Rows 1--3 stay in counts on purpose: the bus bars keep favoring the
ninth graders until the percents arrive. This is also the opening of the unit, so say why it
matters: Unit~3 ends at chance, and chance is built on the right percent.
\end{skillbox}

\boxguard
\begin{skillbox}[Guided Notes \& Practice --- I do, then we do (45 min)]{frost}
\textbf{This is the whole lesson.} There is no group activity, and there is no independent
practice set on this page: \textbf{the homework is the individual practice}, assigned at the
close and done outside class. The notes are \textbf{one Main Ideas / Notes table}: four
instruction rows (problems 1--11), then the Guided Practice row (12--15). \textbf{The rhythm in
every row is the same}: read the definition aloud, mark up the display, work the first problem of
the grid yourself, then hand the rest to the room. The vocabulary box on page~1 is filled
\emph{as each term is named}. The only blanks are the twelve cells of the table in problem~5;
every other answer is written in a problem's own space.
\begin{multicols}{2}
\textbf{I do / we do, row by row (about 25 min).}

\emph{Row 1, Bar Graphs (problems 1--2).} The warm-up's table, drawn twice. Name
\emph{frequency} and \emph{relative frequency}, label the two graphs, and show that the shape is
the same and only the axis changed. Work problem~1 yourself ($94$ students). The class does 2
($80/200 = 40\%$, and the right-hand bar confirms it). Five minutes.

\emph{Row 2, Two Kinds of Bars (problems 3--4).} The same survey split by grade. Work problem~3
yourself (Car, a gap of $16$). \textbf{Problem~4 is a reading trap}: the numbers beside the
segmented bars are running totals, so the twelfth graders' Car segment is $72 - 40 = 32$, not
$72$. Say aloud that both graphs are in counts and that the bars are different heights because
the grades are different sizes; do not go further. Six minutes.

\emph{Row 3, Two-Way Tables (problems 5--7).} Read the summary aloud, then build the table with
the class (problem~5); the Walk column is the subtraction the summary leaves out. Name \emph{joint
count} and \emph{marginal frequency} on the finished table and work problem~6 yourself ($8$, a
joint count). The class does 7 ($80$, a marginal frequency). Six minutes.

\emph{Row 4, Which Whole? (problems 8--11) --- the crux.} Read the two definitions and the bold
sentence aloud. Work problem~8 live ($94/200 = 47\%$). Then \textbf{re-take the hook vote at
problem~9 before anyone divides}. Then compute: $54/120 = 45\%$ of ninth graders, $40/80 = 50\%$
of twelfth graders. \textbf{Problem~9 is the crux}: the group with the larger count has the
smaller rate. Turn the page to the $100\%$ bars and land the boxed sentence word by word.
\textbf{Problem~10 is the trap}: Dani's $54/94$ is a percent of bus riders ($57\%$ of them are
ninth graders), not of ninth graders ($45\%$ of them ride). Problem~11 is the counterweight:
different counts, the same $40\%$. Eight minutes.

\columnbreak
\textbf{We do (about 20 min).} The Guided Practice row, \emph{Streaming Choices} (problems
12--15) --- a fresh context from the Fairview Community Library, given as a pre-made spreadsheet.
\textbf{Students hold the pen; you hold the questions.} Problem~12 is a marginal relative
frequency ($49/250 = 19.6\%$); problem~13 a conditional one ($25/100 = 25\%$). \textbf{Problem~14
is the trap transferred}: Ty's $25/49$ is a percent of Pulse fans. \textbf{Problem~15 is the crux
transferred}, and it comes last: take a hand count on the flyer before anyone divides; then
$60/150 = 40\%$ against $45/100 = 45\%$.

Circulate during Guided Practice and demand the reason, not the word:
\begin{itemize}[leftmargin=1.1em, nosep]
  \item \textbf{Problem 12:} ``What is your whole? Point to it in the table.'' (Cell E4.)
  \item \textbf{Problem 13:} ``Read me the words after \emph{of}.'' (The patrons under $30$:
        cell E2.)
  \item \textbf{Problem 14:} ``What does the $49$ count?'' (Pulse fans of every age.)
  \item \textbf{Problem 15:} ``The count is bigger. So what would decide it?'' A student who
        sides with the flyer has not left the hook vote.
\end{itemize}
While circulating, pick the one paper to put up in the debrief --- a problem~15 that found both
percents and said the flyer compared counts. \textbf{This circulation is your formative read}
--- the homework is not started in class.
\end{multicols}
\end{skillbox}

\boxguard[24]
\begin{skillbox}[Debrief --- whole class, spoken (5 min)]{frost}
Pens down, papers up. \textbf{This debrief is spoken} --- there is no box at the end of the
notes to fill in, so it lives entirely on the board and in the room.
\begin{multicols}{2}
\begin{enumerate}[leftmargin=1.3em, nosep, itemsep=3pt]
  \item \textbf{The one board item: problem~15.} Put up the paper you picked ($40\%$ of patrons
        $30$ and over; $45\%$ of patrons under $30$) and make a student say the sentence:
        \emph{counts say how many; conditional percents say how likely.} Point back at the
        warm-up's line on the board: \emph{larger top number, smaller percent}.
\end{enumerate}
\columnbreak
\textbf{Check for understanding --- ask it cold, whole class.} \emph{``Thirty of $60$ juniors and
$40$ of $100$ seniors signed up for a field trip. Morgan says seniors are more likely to sign
up. Is she right? Give both percents.''} Hands down, thirty seconds of think time, then cold-call
three students. Correct answer: \textbf{no} --- $30/60 = 50\%$ of juniors and $40/100 = 40\%$ of
seniors signed up, so juniors are more likely even though fewer of them signed up. \emph{This
replaces the exit ticket.}

\textbf{Formative read.} Sort the answers into three piles: \textbf{(a)} said \emph{no} and gave
both conditional percents; \textbf{(b)} said \emph{no} but used the wrong whole (for example,
$30/70$); \textbf{(c)} agreed with Morgan --- counts only. \textbf{If pile (c) runs past a third
of the room, open Lesson~3.2 by re-running row~4 (problems 9--10)}, and say so plainly.

\textbf{Five minutes: item~1 and the cold check only.} Anything else waits for the next warm-up.
\end{multicols}
\end{skillbox}

\boxguard
\begin{skillbox}[Homework --- scored 1 to 5, assigned today, due after two study halls]{goldbox}
Six items on \textbf{pet adoptions at the Maple Street Animal Shelter} ($300$ adoptions, dog or
cat, by adopter age): \textbf{1} complete the two-way table from the summary (the core procedure);
\textbf{2} the percent of all adoptions that were dogs, and what it means (\emph{the interpret
item}, $55\%$); \textbf{3} and \textbf{4} are the \textbf{contrast pair}: the same count of $30$
over two different wholes --- of adopters $55$ and over, $60\%$ took a cat; of cats, $22.2\%$
went to adopters $55$ and over; \textbf{5} the director's claim built on counts --- \emph{the
crux head-on and a justification} ($40\%$ vs.\ $60\%$); \textbf{6} which of two pre-drawn displays
compares the age groups --- \emph{the spiral item}, reaching back to Lesson~2.4's common scale,
and a second justification.

\textbf{Assigned at Close \& Assign and done outside class.} \textbf{Scoring:} 12 points --- 2
per turned-in item and the DeltaMath score scaled to the other 6, which the grading report turns
into the score of 1 to 5. \textbf{Item~5 is where the misconception shows}: a student who agrees
with the director is still comparing counts. \textbf{The second check (items 3--4, which whole)
moves to DeltaMath's Percents Two-Way Tables}: a student who misses the column-percent problems is
dividing by the wrong total.

\textbf{DeltaMath + justifications --- decided here.} The packet pages are authored and are the
default; no Desmos activity asks for the written reasons this lesson exists for. Item~1 is the
two-way-table procedure and items 3--4 are conditional percents with their denominators, which
DeltaMath drills; items 2, 5, and~6 are an interpretation, the crux, and a choice of display, which
it cannot grade --- and DeltaMath has no segmented or percent bar chart topic, so item~6 comes in on
paper. So assign \textbf{DeltaMath} --- Fill in Values, Two-Way Tables (L1) $\cdot$ Fill in Values,
Two-Way Tables (L2) $\cdot$ Interpret Two-Way Tables $\cdot$ Percents Two-Way Tables --- for the
computation (items 1, 3, and~4 stay in the packet as a worked reference), and students
\textbf{turn in packet items 2, 5, and~6}. Still 12 points: 2 per turned-in item, the DeltaMath
score scaled to the other 6. Say so aloud at Close \& Assign.

\textbf{Preview of the next lesson.} Lesson~3.2, Probability Models: a relative frequency is the
first step toward a probability.

\textbf{Connections \& Big Ideas.} \textit{Unit 3 core idea:} data sorted into groups give chances,
and every chance is a part over a whole. \textit{Standards covered:} PS.DS.4a (bar graphs of one
categorical variable), PS.DS.4b (side-by-side and segmented bar graphs comparing groups),
PS.DS.4c (generate and interpret a two-way table), PS.DS.4d (marginal, relative, and conditional
frequencies in context). \textit{Spirals forward to:} Lessons 3.2 (probability models) and 3.4
(conditional probability and independence). \textit{Reaches back to:} Lesson~2.4's rule that
groups are compared on a common scale, and prior-course percent work.
\end{skillbox}

\boxguard
\begin{skillbox}[Watch For (while circulating)]{redbox}
\begin{itemize}[leftmargin=1.2em, nosep]
  \item \textbf{Comparing raw counts across groups of different sizes} (the hook, problems 9, 11,
        and~15; HW~5 --- the target misconception). Probe: \emph{``How many ninth graders did we
        ask? How many twelfth graders? So what would a fair comparison use?''}
  \item \textbf{The wrong whole} (problems 10 and~14; HW~3--4 --- the second misconception):
        $54/94$ reported as a percent of ninth graders. Probe: \emph{``Read me the words after
        \emph{of}. Where is that total in the table?''}
  \item \textbf{A running total read as a segment} (problem~4): $72$ for the twelfth graders' Car
        riders. Probe: \emph{``Where does the Car segment start, and where does it stop?''}
  \item \textbf{A joint count called a marginal frequency} (problems 6--7): Probe: \emph{``Is that
        number inside the table or in a Total?''}
  \item \textbf{A percent with no group} (problems 2, 12, 13; HW~2): ``$40\%$ come by car.''
        Probe: \emph{``Forty percent of whom?''}
  \item \textbf{Picking the count display} (HW~6): Probe: \emph{``If the shelter had twice as many
        adopters in one age group, what would happen to that bar?''}
\end{itemize}
Cold-call prompts: \emph{``Give me a group at this school where a bigger count could hide a
smaller percent.''} \quad \emph{``When is a raw count the number you actually want?''}
\end{skillbox}

\boxguard
\begin{skillbox}[Close \& Assign (5 min)]{goldbox}
\textbf{The homework is assigned here and done outside class --- it is not started in class.}
Hand the launch: six items on pet adoptions, scored 1 to 5, due at the start of the first class
after two study halls. Say the form: \textbf{DeltaMath} for the two-way tables and their
percents, and \textbf{packet items 2, 5, and~6 turned in}. Point out that the \emph{Keep in Mind} box on the
cover and rows 3--4 of the notes are the pages to flip back to. \textbf{Tell students to do item~5
first if they are short on time} --- it shows whether row~4 landed --- and that items 3 and~4 are
a pair: read both before answering either. Then close on what changed: \emph{``You walked in
reading counts. You walk out asking, of whom? Counts say how many; conditional percents say how
likely.''} \textbf{Preview:} Lesson~3.2, Probability Models --- the relative frequency you found
today is the first step toward a probability.
\end{skillbox}

% ── Teacher notes — one per component, in packet order (three) ──────────────
% Teacher-only prose lives HERE, never in a _key. No note for the debrief — it is fully
% specified in its own box above.
\begin{teachernote}[Warm-Up]
Five minutes, silent, timed. This is the first page of Unit~3, so open with one sentence: the
unit is about data sorted into groups and the chances those data give us. The old unit launch is
folded into this warm-up and the hook; there is no separate launch day. Items 1 and~3 should be
quick. \textbf{Item~2's last question is the one to read}: students who write that the larger
top number is the larger percent have walked in with the target misconception --- note two of
them and cold-call them at problem~9. The warm-up prints no grade sizes in context, on purpose,
so the hook stays open. Write \emph{larger top number, smaller percent} on the board and
\textbf{leave it up all period}; the debrief points back at it.
\end{teachernote}

\begin{teachernote}[Guided Notes \& Practice]
\textbf{This one document is the lesson --- pace it 25 / 20, then 5 for the debrief and 5 to
assign the homework.} It is one table, and every row runs the same way: you read the definition,
mark up the display, and work the first problem of the grid (1, 3, 6, 8); the class works the
rest with the pen in their hand. Rows 1 and~2 must move --- five and six minutes. All four rows
are one survey, so students are reading one data set four ways, not learning four data sets.

\textbf{Spend the time in row~4, \emph{Which Whole?}}, which carries the misconception. Rows 1--3
are in counts on purpose. \textbf{Re-take the hook vote at problem~9 before anyone divides.} If
students divide first, the vote is empty; if they commit and \emph{then} find $45\%$ against
$50\%$, the rule arrives as the answer to a question they already have. The $100\%$ bars are
printed after problem~9, at the top of the next page, so the picture confirms the computation
rather than replacing it. \textbf{Problem~10 (the wrong whole) is not a side trip}: it is the
second way the same mistake appears, and Guided Practice~14 and homework items 3--4 ask for it
again. Problem~11 keeps the rule honest: sometimes the counts differ and the likelihood does not.

\textbf{Guided Practice (12--15) is where the pen changes hands.} \emph{Streaming Choices} is a
fresh context given as a spreadsheet, so students must find each whole in the table themselves.
Problem~15, last, is the crux transferred: take a hand count on the flyer before anyone divides.
If time is short, protect 14 and 15 --- problem~12 can be said aloud. Creating graphs in software
is not worked here; the charts are pre-made output, and homework item~6 is the last place
students read one in this lesson.

\textbf{There is no independent practice set on this page, and that is deliberate --- the
homework is the individual practice}, assigned at the close and done outside class. Your
formative read is the Guided Practice circulation plus the cold check; if the counts-only pile is
more than a third of the room, \textbf{open Lesson~3.2 by re-running problems 9--10}, and say so
plainly.
\end{teachernote}

\begin{teachernote}[Homework]
\textbf{Scored 1 to 5 from 12 points, due at the start of the first class after two study halls;
assigned at the close, not started in class.} Items 2, 5, and~6 are the turned-in rubric items,
2 points each (see \texttt{rubrics/}). \textbf{Item~5 is the one to read
first.} Sort it into three piles --- found both percents and named the different group sizes
(ready); found both percents but gave no reason (ready, needs the language); agreed with the
director or used the wrong whole (the misconception survived). If more than a third land in the
last pile, open Lesson~3.2 with problems 9--10 before anything else.

\textbf{DeltaMath + justifications is the call for 3.1.} DeltaMath (exact topic names): Fill in
Values, Two-Way Tables (L1) $\cdot$ Fill in Values, Two-Way Tables (L2) $\cdot$ Interpret Two-Way
Tables $\cdot$ Percents Two-Way Tables. Students turn in packet items 2, 5, and~6; items 1, 3,
and~4 stay in the packet, unboxed, as a worked reference. Score the three turned-in items as above
(item~5 first, then 6, then 2); scale the DeltaMath score to the other 6 points. The contrast-pair
check on items 3--4 (the same $30$ over two wholes) moves to DeltaMath's Percents Two-Way Tables:
a student who misses the column-percent problems is dividing by the wrong total and has not yet
asked \emph{of whom?} Item~6 stays on paper because DeltaMath has no segmented or percent bar
chart topic. Open the next class by putting two anonymous item~5 answers from different piles
side by side on the board.
\end{teachernote}
\end{document}
